\section{Receiver messages and concavification}
\label{sec:messages96}
The distinction between a device label and a receiver message is
operational, not merely terminological.  A device label may select a
fresh source even when an intermediate receiver result may not.  We
first give the exact postponement statement, then characterize when
this additional communication has value in an arbitrary finite
experiment.  The old receiver is large enough to retain the complete
initial Schmidt support throughout this section.  Only the fresh
reference is dimension constrained.

\subsection{Postponement without a receiver message}
Let $R$ be the initial receiver, $I_2$ the second input, and $F_y$ the
fresh reference after the first device label $y$.  The prescribed fresh
state $\sigma_y$ on $I_2\otimes F_y$ is independent of a receiver outcome
$h$.  A normal instrument $\mathcal K_{h|y}:\mathcal T(R)\to
\mathcal T(O_{yh})$ acts on the old receiver.  Write
$M_{j|yhz}$ for the final decision POVM on $O_{yh}\otimes F_y$.
The record $(y,h,z)$ is classical; retaining $h$ for the final readout
is permitted even though it is not sent to the source.

\begin{lemma}[No-message postponement]\label{lem:postponement96}
For every such instrument and readout, the operators
\begin{equation}\label{eq:heisenberg96}
 G_{j|yz}=\sum_h
       (\mathcal K_{h|y}^{*}\otimes\operatorname{Id}_{F_y})(M_{j|yhz})
\end{equation}
form a POVM on $R\otimes F_y$.  Replacing the instrument and its record
by this final POVM preserves every decision probability, under each
hypothesis separately, for every input state and every second device.
The statement includes countable normal instruments and their
finite-dimensional tester closure.  It preserves any fresh-reference
dimension bound.  It preserves an old-receiver bound $k$ whenever
$\dim R\le k$, or after compression to a common input support of
dimension at most $k$.
\end{lemma}
\begin{proof}
Every summand in \eqref{eq:heisenberg96} is positive.  Since the sum of
the instrument maps is trace preserving,
\[
 \sum_jG_{j|yz}=\sum_h\mathcal K_{h|y}^{*}(I_{O_{yh}})
                                      \otimes I_{F_y}=I_R\otimes I_{F_y}.
\]
Let $X_{\theta,y}$ be an arbitrary unnormalized first receiver block
and let $Z_{\theta,y,z}$ be the fresh reference block obtained by
applying the second device to $\sigma_y$.  Conditional independence
at the preparation cut gives
\[
 \begin{split}
 &\sum_h\tr\bigl[M_{j|yhz}
        (\mathcal K_{h|y}(X_{\theta,y})\otimes Z_{\theta,y,z})\bigr]\\
 &\hspace{35mm}=\tr[G_{j|yz}(X_{\theta,y}\otimes Z_{\theta,y,z})].
 \end{split}
\]
This is the adjoint identity, with no division by a history
probability.  It proves the assertion also at null histories.
For countably many outcomes the sums of positive effects increase
weakly to bounded effects, and the corresponding probabilities
converge by Lemma~\ref{lem:countable91}.  Equivalently retain a remainder
branch and pass to its vanishing trace tail.  Continuity then covers
norm limits of finite testers.  No new quantum system was introduced;
the old system retained by the replacement is precisely $R$, and the
fresh reference remains $F_y$.  Compression to a common support and
extension of the POVM by an arbitrary decision on its complement give
the last assertion.
\end{proof}
A constraint only on the smaller \emph{post-instrument} system is not
sufficient for this resource statement: postponing a compression can
retain a larger $R$.  We do not apply the lemma to bypass such a cut.
When $\sigma_{yh}$ depends on $h$, there need not be one common
$Z_{\theta,y,z}$, so the displayed factorization is unavailable.
First-label dependence $\sigma_y$ causes no such difficulty.

\subsection{An exact criterion for arbitrary finite experiments}
Retain the experiment and the positive payoff matrices
\eqref{eq:gameblocks91}.  Fix a pure first acquisition $|C_0\rangle$
with $C_0:I_1\to R$, $C_0^*C_0=\rho$, and compress $R$ to the range
of $C_0$.  The initial state is held fixed, not replaced by an initially
announced mixture.  Let $1\le\ell\le d_2$ and write
\[
 g_{y,\ell}(a)=
  \max_{b\in\mathcal D_{d_2},\ \operatorname{rank}b\le\ell}\Phi_y(a,b),
 \qquad
 c_{y,\ell}(\rho)=
  \max_{\substack{\rho=\sum_hw_ha_h\,,\ w_h\ge0\\a_h\in\mathcal D_{d_1}}}
                                      \sum_hw_hg_{y,\ell}(a_h).
\]
Here $\Phi_y$ is the final-POVM optimization in
\eqref{eq:leafvalue91}.  The variables $w_h=\tr A_h$ are
\emph{Gram weights}; the physical first-record probability is
$w_h\tr(a_h(E_y^{\theta,1})^{\mathsf T})$.  The two are not
interchangeable.  Both old receivers and their intermediates may be
retained without a bound below $d_1$.  Write
$P_{\mathrm{label},\ell}(\rho)$ and
$P_{\mathrm{msg},\ell}(\rho)$ for the fixed-initial optima when fresh
preparation may use only $y$, or may also use the receiver message.

\begin{theorem}[Uniform message criterion]\label{thm:messagecriterion96}
For every finite experiment, every prescribed pure initial Gram
$\rho$, and every fresh-reference bound $\ell$,
\begin{equation}\label{eq:messagecriterion96}
 \begin{split}
 P_{\mathrm{label},\ell}(\rho)&=\sum_yg_{y,\ell}(\rho),\\
 P_{\mathrm{msg},\ell}(\rho)&=\sum_yc_{y,\ell}(\rho),\\
 \Delta_\ell(\rho)&=
       \sum_y\bigl[c_{y,\ell}(\rho)-g_{y,\ell}(\rho)\bigr]\ge0.
 \end{split}
\end{equation}
Both optima are attained; for the second, at most
$(\operatorname{rank}\rho)^2$ nonzero receiver outcomes after each
first label suffice.  The following statements are equivalent:

\smallskip
\noindent (i) $\Delta_\ell(\rho)=0$ for every prescribed pure first
acquisition $\rho$;

\noindent (ii) every $g_{y,\ell}$ is concave on $\mathcal D_{d_1}$;

\noindent (iii) no protocol with a binary receiver instrument after
one first label, and no receiver instrument after the other labels,
can strictly outperform the optimal label-only value at its own
prescribed initial acquisition.

\smallskip
If messages help somewhere, two outcomes after a single first label
therefore suffice to witness a strict improvement at \emph{some}
initial acquisition.  This last assertion does not fix in advance the
initial acquisition at which the witness occurs.
\end{theorem}
\begin{proof}
Square roots, payoff contractions and the compact final-POVM
optimization make $\Phi_y$ continuous.  The set of rank-at-most-$\ell$
densities is compact, so $g_{y,\ell}$ is continuous and its maximum is
attained.  A mixed fresh preparation with reference dimension $\ell$
may be decomposed into pure vectors on that same space.  With no
receiver message, their mixture is independent of the old record.
Linearity shows that one of these fresh vectors does at least as well
as the mixture for a fixed continuation.  Its Gram has rank at most
$\ell$; conversely every such Gram has a purification with reference
dimension $\ell$.  Lemma~\ref{lem:postponement96} now proves the first
formula without disallowing old quantum information in the final
readout.

With messages, resolve receiver Kraus indices and fresh pure-state
mixtures into the public record.  Every fresh pure coefficient map has
rank at most $\ell$ on its existing reference; refining the receiver
instrument keeps its output space unchanged.  The resulting branch
Grams have the common sum $\rho$ for every $y$, and the best branch
value is $w_hg_{y,\ell}(a_h)$.  Conversely
Lemma~\ref{lem:resetnormal91} implements every finite such decomposition
on the fixed Schmidt support; the fresh factors have rank at most
$\ell$.  The graph of $g_{y,\ell}$ on the density matrices supported
in $\operatorname{supp}\rho$ is compact.  Its convex hull gives an
attaining roof.  If an attaining ensemble has more than
$s^2$ atoms, $s=\operatorname{rank}\rho$, its density matrices are
linearly dependent in a real vector space of dimension $s^2$.
Vary their weights in both directions along this dependence.  An
objective-changing direction would contradict optimality; otherwise
move until a weight vanishes.  Iterating proves the atom bound.
Countable outcomes follow by positive tails.  These arguments prove
the second formula and its attainment.

Each difference in the last line of \eqref{eq:messagecriterion96} is
nonnegative.  If (i) holds, each difference vanishes at every $\rho$;
thus $g_{y,\ell}=c_{y,\ell}$ and (ii) follows.  Conversely concavity
makes each roof equal to its defining function, proving (i).
If (ii) fails, by the definition of concavity there are a label $y_0$,
states $a,b$, and $0<\lambda<1$ such that, for
$\rho=\lambda a+(1-\lambda)b$,
\[
 \lambda g_{y_0,\ell}(a)+(1-\lambda)g_{y_0,\ell}(b)
                                      >g_{y_0,\ell}(\rho).
\]
Realize this two-atom instrument after $y_0$, with its optimizing fresh
continuations.  At every other first label use the optimizing
label-only continuation at $\rho$.  The strict displayed excess is
exactly its excess over the \emph{optimal} label-only value at $\rho$.
This disproves (iii).  Finally (i) implies (iii) by class inclusion.
\end{proof}
The criterion applies to fixed-pair as well as ensemble problems, but
it does not assert that a particular experiment has a positive gap.
It classifies uniform absence of message advantage in the stated
reset interface, not every meaning of bounded coherent memory.
The inclusion diagram at a fixed admissible initial acquisition is
\[
 \begin{array}{ccccc}
 \mathfrak C_\ell(\rho)&\subseteq&\mathfrak R_{d_1,\ell}(\rho)
                                      &\supseteq&\mathfrak L_\ell(\rho),\\
 &&\mathfrak R_{d_1,\ell}(\rho)\subseteq\mathfrak A(\rho).&&
 \end{array}
\]
where $\mathfrak C$ keeps only classical old information,
$\mathfrak L$ allows label feedback and full old quantum storage but
no receiver-to-source message, and $\mathfrak A$ allows coherent
acquisition.  No inclusion between $\mathfrak C$ and $\mathfrak L$
is used.  At a fixed Gram, a message gap is compatible with a globally
optimized no-message strategy using a different first acquisition.

\begin{corollary}[Stability of the message gap]\label{cor:messageperturb96}
Keep the labels, reward table, prior law, and both optimization classes
fixed.  For every latent index $\theta$, suppose the entire slot-$i$
measurement channel changes by at most $\delta_i$ in \emph{unhalved}
diamond norm, $i=1,2$.  The two optima in
\eqref{eq:messagecriterion96} each change by at most
$(\delta_1+\delta_2)/2$, and
\[
       |\widetilde\Delta_\ell(\rho)-\Delta_\ell(\rho)|
                                             \le\delta_1+\delta_2.
\]
In particular a larger original gap remains strict under these
perturbations.  The latent device is still selected once.
\end{corollary}
\begin{proof}
For a fixed protocol, replace the two channel slots in turn.
Diamond-norm control, including every reference, and trace-norm
contractivity bound the change of its final normalized state by
$\delta_1+\delta_2$.  The difference has trace zero, so a reward effect
between zero and the identity changes its expectation by at most
one half of that norm.  Average over the fixed prior and take the
supremum in each unchanged class.  The triangle inequality gives
the gap bound.  The argument compares specified valid reset protocols;
it does not infer physical reset independence from the observed score.
\end{proof}
